\begin{center}
Problema I

Heverton e a Fila 

{\tiny by André da Cunha Ribeiro}

Timelimit: $1$	
\end{center}			

Heverton é um coordenador de curso muito dedicado do curso de computação do NuCAT. Os alunos estão enfrentando um grande problema na entrada do NuCAT, todos querem entrar no mesmo horário gerando assim um grande tumulto. Heverton reuniu com os cavalheiros do NDE para tentar resolver o problema. Depois de muita discussão o NDE resolveu ordenar os alunos pelo número de matricula, e organizando a entrada pela ordem decrescente. Mas Heverton precisa de sua ajuda para desenvolver o programa pois ele está muito ocupado com seus trabalhos.

\vspace{0.3cm}
\textbf{Entrada}
Para cada dia será lido um inteiro $N$, $1 \leq  N  \leq  10^{4}$, que representa o número de alunos disponíveis na entrada. Depois será lido os $N$ números de matriculas.

\vspace{0.3cm}
\textbf{Saída}
Para cada dia, imprimir os  número de  matriculas ordenados em ordem decrescente.

\begin{table}[h]
    \centering
    \begin{tabular}{|p{7cm}|p{7cm}|}
        \hline
        \begin{center}
            \vspace{-0.3cm}
            \textbf{Exemplos de Entrada}
            \vspace{-0.3cm}
        \end{center} 
         &  
        \begin{center}
            \vspace{-0.3cm}
            \textbf{Exemplos de Saída}
            \vspace{-0.3cm}
        \end{center} \\ \hline
         \texttt{5} & \texttt{5 4 3 2 1} \\  
         \texttt{1 2 3 4 5} & \texttt{} \\
         \hline
    \end{tabular}
    \caption{Exemplos de entradas e saídas}
    \label{tab:inputI2}
\end{table}

\newpage

\begin{center}
\noindent
\lstinputlisting{abobora_23/I/I3.cpp}
\end{center}
